% !TEX root = ../../main.tex

\section{Risk and Return Fundamentals}

\subsection{Overview}

In finance, the fundamental relationship between risk and return is central to all investment decisions. Risk is the uncertainty of outcomes, while return is the gain (or loss) from an investment. Investors demand compensation—a risk premium—for bearing risk. The question is: how do we measure this trade-off accurately?

\subsection{Measuring Risk}

\subsubsection{Standard Deviation and Total Risk}

\textbf{Standard Deviation ($\sigma$):} Measures the dispersion of a dataset relative to its mean. In finance, it represents \textbf{total risk}—the combination of systematic risk (market-wide) and idiosyncratic risk (firm-specific).

For a return series, standard deviation quantifies how much returns deviate from the average, on average. Higher standard deviation indicates more volatility and unpredictability.

\subsubsection{Risk Decomposition}

Total risk can be decomposed into two components:
\begin{itemize}
    \item \textbf{Systematic Risk ($\beta$):} Market-wide risk that cannot be eliminated through diversification. Examples: economic cycles, interest rate changes, inflation.
    \item \textbf{Idiosyncratic Risk ($\epsilon$):} Firm-specific risk that can be diversified away through holding multiple assets. Examples: management changes, product recalls, company-specific litigation.
\end{itemize}

\subsection{Risk-Adjusted Performance Measures}

Performance evaluation requires more than just looking at returns. A portfolio earning 10\% with 20\% volatility is more impressive than one earning 10\% with 50\% volatility. \textbf{Risk-adjusted performance measures} quantify how efficiently a manager converts risk into returns.

\subsubsection{Sharpe Ratio}

The Sharpe Ratio measures the excess return per unit of total risk:
\begin{equation}
    \text{SR} = \frac{\mu_p - r_f}{\sigma_p}
\end{equation}

where:
\begin{itemize}
    \item $\mu_p$ = expected return of the portfolio
    \item $r_f$ = risk-free rate (e.g., Treasury yield)
    \item $\sigma_p$ = standard deviation (volatility) of portfolio returns
\end{itemize}

The numerator $(\mu_p - r_f)$ is called the \textbf{excess return} or \textbf{risk premium}—the additional return earned for taking on risk beyond the risk-free asset.

\paragraph{Components Explained}
\begin{enumerate}
    \item \textbf{Expected Return ($\mu_p$):} The average return you expect from the portfolio over a given period. For historical data, use the mean of past returns. For forward-looking analysis, use analyst forecasts or model-based projections.
    
    \item \textbf{Risk-Free Rate ($r_f$):} The return on a zero-risk investment, typically a government bond (e.g., US Treasury). This is the minimum return an investor can earn with certainty. The choice of maturity should match the investment horizon.
    
    \item \textbf{Volatility ($\sigma_p$):} The standard deviation of returns, measuring both upside and downside variability. A higher volatility means returns are more unpredictable.
\end{enumerate}

\paragraph{How to Interpret the Sharpe Ratio}
\begin{itemize}
    \item \textbf{Positive SR $> 1$:} The portfolio earns significantly more return per unit of risk—generally considered attractive. A SR of 1 means for every unit of risk, you earn 1 unit of excess return.
    
    \item \textbf{SR between 0 and 1:} The portfolio adds value, but modestly. Risk is being compensated but not generously.
    
    \item \textbf{SR $\leq 0$:} The portfolio underperforms the risk-free rate on a risk-adjusted basis. You'd be better off holding Treasury bills.
    
    \item \textbf{Negative SR:} The portfolio returns are below the risk-free rate, which is even worse (you take on volatility for negative excess returns).
    
    \item \textbf{Comparing portfolios:} A portfolio with SR = 1.5 is more efficient than one with SR = 0.8. The higher ratio indicates better risk-adjusted performance.
\end{itemize}

\paragraph{Where and How It Is Used}
\begin{enumerate}
    \item \textbf{Portfolio Comparison:} Investors use Sharpe Ratios to rank funds and strategies on a level playing field. Two funds with different risk levels can be fairly compared via their SR.
    
    \item \textbf{Asset Manager Performance Evaluation:} Trustees and investors assess whether fund managers are earning their fees by comparing a fund's SR against benchmarks and peers. A fund with higher SR relative to its category suggests superior stock-picking or risk management.
    
    \item \textbf{Portfolio Optimization:} The Sharpe Ratio is the objective function in Modern Portfolio Theory. Markowitz's efficient frontier is constructed by maximizing the Sharpe Ratio for a given level of risk, identifying the portfolio with the best risk-return trade-off.
    
    \item \textbf{Risk Control and Rebalancing:} Portfolio managers monitor the SR over time. If it deteriorates, it signals that either returns have fallen or risk has increased disproportionately, prompting a rebalancing decision.
    
    \item \textbf{Hedge Fund and Systematic Strategy Evaluation:} Especially in quantitative and alternative asset management, the SR is the primary yardstick. A 1.0+ SR is often considered the bar for professional-grade returns.
    
    \item \textbf{Performance Attribution:} If a strategy's SR degrades, analysts decompose it into changes in return vs. increases in volatility to diagnose the cause.
\end{enumerate}

\paragraph{Limitations and Considerations}
\begin{itemize}
    \item \textbf{Assumes normal distributions:} Sharpe assumes returns are normally distributed. In reality, market returns often have fat tails and skewness, especially during crises.
    
    \item \textbf{Penalizes all volatility equally:} It treats upside and downside volatility the same. Investors generally prefer upside surprises, so the Sortino Ratio (penalizing only downside deviation) may be more appropriate.
    
    \item \textbf{Backward-looking:} Historical Sharpe Ratio does not guarantee future performance.
    
    \item \textbf{Time period sensitivity:} Changing the measurement period (daily, monthly, annual) can significantly alter the SR due to different volatility regimes.
    
    \item \textbf{Risk-free rate assumption:} The choice of $r_f$ matters. Using a 10-year Treasury when evaluating a 1-year strategy creates a mismatch.
\end{itemize}

\subsubsection{Sortino Ratio}

Unlike the Sharpe Ratio, the Sortino Ratio only penalizes "bad" volatility by using downward deviation ($\sigma_d$), meaning it doesn't punish an investor for large positive "surprises." This makes it more suitable for assessing strategies where upside volatility is desirable or for investors who care primarily about downside protection.

\section{Understanding Market Risk: Systematic vs. Idiosyncratic}

\subsection{Risk Decomposition in the Market}

Earlier, we noted that total risk has two components. Understanding this distinction is crucial for portfolio management and pricing models like CAPM.

\begin{itemize}
    \item \textbf{Systematic Risk:} Cannot be eliminated through diversification. Driven by economy-wide factors. Investors must be compensated for bearing it.
    
    \item \textbf{Idiosyncratic Risk:} Can be eliminated through diversification. Driven by company-specific factors. Rational investors should not demand a premium for bearing it (since they can avoid it).
\end{itemize}

This distinction motivates the \textbf{Capital Asset Pricing Model}, which relates expected return only to systematic risk.

\section{The Capital Asset Pricing Model (CAPM)}

\subsection{The CAPM Framework}

The Capital Asset Pricing Model describes the relationship between an asset's systematic risk and its expected return:
\begin{equation}
    E[R_i] = R_f + \beta_i (E[R_m] - R_f) + \alpha_i
\end{equation}

This equation has several interpretations:
\begin{itemize}
    \item The expected return on any asset equals the risk-free rate plus a risk premium.
    \item The risk premium is proportional to the asset's systematic risk (beta) and the market risk premium.
    \item Alpha represents the return above or below what CAPM predicts.
\end{itemize}

\subsection{Beta ($\beta$): Systematic Risk}

Beta measures the sensitivity of an asset's returns to the overall market. It is the slope of the security characteristic line when regressing asset returns against market returns:
\begin{equation}
    \beta_i = \frac{\text{Cov}(R_i, R_m)}{\text{Var}(R_m)} = \rho_{im} \frac{\sigma_i}{\sigma_m}
\end{equation}

\paragraph{Interpreting Beta}
\begin{itemize}
    \item \textbf{$\beta = 1$:} The asset moves perfectly with the market. Neither more nor less risky than the market.
    
    \item \textbf{$\beta > 1$:} The asset is more volatile than the market (amplifies market moves). Riskier and demands higher expected return.
    
    \item \textbf{$0 < \beta < 1$:} The asset is less volatile than the market (dampens market moves). Less risky.
    
    \item \textbf{$\beta < 0$ or $\beta = 0$:} The asset is uncorrelated with the market (negative correlation or zero correlation). Valuable for diversification.
\end{itemize}

\subsection{Alpha ($\alpha$): Abnormal Returns}

Alpha is the intercept of the regression. It represents the "skill" or "edge"—the return above or below what CAPM predicts:
\begin{itemize}
    \item \textbf{$\alpha > 0$:} The asset outperforms its predicted return. The manager adds value.
    \item \textbf{$\alpha = 0$:} The asset performs exactly as CAPM predicts. Consistent with market efficiency.
    \item \textbf{$\alpha < 0$:} The asset underperforms its predicted return. The manager destroys value.
\end{itemize}

In a perfectly efficient market (under the Efficient Market Hypothesis), expected alpha is zero. Persistent positive alpha is usually attributed to:
\begin{itemize}
    \item Information asymmetry (better data or research)
    \item Structural advantages (lower transaction costs, latency, or fees)
    \item Superior analytical logic or predictive models
\end{itemize}

\subsection{The Security Market Line (SML)}

The SML is a graphical representation of CAPM, plotting expected return against beta:
\begin{equation}
    E[R_i] = R_f + \beta_i (E[R_m] - R_f)
\end{equation}

Assets lying above the SML have positive alpha (underpriced), while those below have negative alpha (overpriced). The slope of the line is the market risk premium $(E[R_m] - R_f)$.

\section{Portfolio Risk Management}

\subsection{Diversification: The Power of Correlation}

The fundamental principle of portfolio management is that combining assets with imperfect correlation reduces risk. This is not intuitive in everyday life but is one of the most powerful findings in finance.

\subsubsection{Two-Asset Case}

Consider two assets with returns $R_A$ and $R_B$, volatilities $\sigma_A$ and $\sigma_B$, and correlation $\rho_{AB}$. A portfolio with weights $w_A$ and $w_B$ has volatility:

\begin{equation}
    \sigma_p = \sqrt{w_A^2 \sigma_A^2 + w_B^2 \sigma_B^2 + 2 w_A w_B \sigma_A \sigma_B \rho_{AB}}
\end{equation}

\paragraph{Key Insights}
\begin{itemize}
    \item If $\rho_{AB} = 1$ (perfect positive correlation), the portfolio risk is the weighted average of individual risks: $\sigma_p = w_A \sigma_A + w_B \sigma_B$. No diversification benefit.
    
    \item If $\rho_{AB} < 1$ (imperfect correlation), the portfolio risk is strictly less than the weighted average. Diversification reduces risk.
    
    \item If $\rho_{AB} = -1$ (perfect negative correlation), you can construct a zero-risk portfolio (if variances match appropriately).
\end{itemize}

The diversification benefit is maximized when assets are uncorrelated or negatively correlated.

\subsubsection{Multiple Assets: Matrix Notation}

For a portfolio of $N$ assets with a weight vector $\mathbf{w}$ and covariance matrix $\mathbf{\Sigma}$:
\begin{equation}
    \sigma_p^2 = \mathbf{w}^T \mathbf{\Sigma} \mathbf{w}
\end{equation}

where the covariance matrix element $\Sigma_{ij} = \sigma_i \sigma_j \rho_{ij}$ captures both individual volatilities and correlations.

This formulation is the foundation of portfolio optimization algorithms that find the \textbf{efficient frontier}—the set of portfolios with maximum return for a given risk level, or minimum risk for a given return level.

\section{Risk Measurement and Quantification Theory}

\subsection{What Makes a Risk Measure "Coherent"?}

Risk managers use various metrics to quantify exposure. However, not all risk measures are created equal. A ``good'' risk measure should satisfy certain mathematical properties that align with intuition. These are captured by the concept of \textbf{coherence}.

\subsection{Coherent Risk Measures}

A risk measure $\rho$ is considered coherent if it satisfies four axioms defined by Artzner et al. (1999):

\begin{enumerate}
    \item \textbf{Monotonicity:} If portfolio $P_1$ always performs at least as well as portfolio $P_2$, then its risk should not exceed $P_2$'s risk. Formally: If $P_1 \leq P_2$ almost surely, then $\rho(P_2) \leq \rho(P_1)$.
    
    \item \textbf{Subadditivity:} The risk of a combined portfolio should not exceed the sum of individual risks. Formally: $\rho(P_1 + P_2) \leq \rho(P_1) + \rho(P_2)$. This reflects the principle that diversification reduces (or maintains) risk.
    
    \item \textbf{Positive Homogeneity:} Scaling a position scales its risk proportionally. Formally: For $\lambda \geq 0$, $\rho(\lambda P) = \lambda \rho(P)$. Doubling a position doubles its risk.
    
    \item \textbf{Translational Invariance:} Adding risk-free cash reduces risk by exactly that amount. Formally: For $m \in \mathbb{R}$, $\rho(P + m) = \rho(P) - m$. This makes intuitive sense: holding cash is the safest investment.
\end{enumerate}

\subsection{Value at Risk (VaR): Definition and Calculation}

\textbf{Value at Risk (VaR)} is one of the most widely used risk metrics in finance. It answers the question: ``How much can I lose with a given confidence level over a given time horizon?''

\subsubsection{Definition}

VaR at confidence level $\alpha$ and time horizon $T$ is:
\begin{equation}
    \text{VaR}_\alpha(T) = -\text{quantile}_\alpha(\text{P\&L})
\end{equation}

For example, a 95\% one-day VaR of \$1 million means there is only a 5\% chance of losing more than \$1 million in a single day.

\subsubsection{Calculation Methods}

\begin{enumerate}
    \item \textbf{Parametric (Variance-Covariance):} Assumes returns are normally distributed. Simple but fails when tails are fat.
    \item \textbf{Historical Simulation:} Uses historical returns as a distribution. Captures actual tail behavior but requires sufficient historical data.
    \item \textbf{Monte Carlo:} Simulates future scenarios using a stochastic model. Flexible but computationally intensive.
\end{enumerate}

\subsection{The Critical Flaw: Why VaR Is Not Coherent}

\textbf{VaR is NOT a coherent risk measure.} It violates the subadditivity axiom, meaning diversification does not always reduce VaR. This is a serious flaw for risk management.

\paragraph{Demonstration: The Digital Bonds Example}

Imagine two independent digital bonds that default with a 4\% probability each (lose 100\% if they default). Set a 95\% confidence level for VaR.

\begin{itemize}
    \item \textbf{Bond A alone:} The probability of default is 4\%, which is below the 5\% threshold. VaR = 0.
    \item \textbf{Bond B alone:} Similarly, VaR = 0.
    \item \textbf{Both bonds (portfolio):} The probability of at least one defaulting is $1 - (0.96 \times 0.96) = 7.8\%$, which exceeds the 5\% threshold. VaR $> 0$.
\end{itemize}

Thus: $\text{VaR}(A + B) > \text{VaR}(A) + \text{VaR}(B)$

This violates subadditivity! The model incorrectly suggests that combining two safe positions creates risk.

\subsection{Expected Shortfall (CVaR): The Coherent Alternative}

\textbf{Expected Shortfall (ES)}, also called Conditional Value at Risk (CVaR), addresses VaR's shortcomings. It measures the \textbf{average loss in the tail beyond the VaR threshold}:

\begin{equation}
    \text{ES}_\alpha = \mathbb{E}[\text{Loss} \mid \text{Loss} > \text{VaR}_\alpha]
\end{equation}

\paragraph{Key Advantages}
\begin{itemize}
    \item \textbf{Coherent:} ES satisfies all four axioms, making it theoretically sound.
    \item \textbf{Captures tail risk:} It considers the severity of losses beyond VaR, not just their probability.
    \item \textbf{Incentivizes proper risk management:} Unlike VaR, ES penalizes strategies that concentrate risk in the tail.
\end{itemize}

\section{Advanced Risk Topics}

\subsection{Derivatives and Counterparty Risk}

\subsubsection{XVA (Valuation Adjustments)}

When trading derivatives, the agreed-upon price is typically the ``mid-market'' price—the midpoint between bid and ask in a perfectly liquid market. However, reality is more complex. Several costs and risks are not captured in the mid-market price, collectively called XVA:

\begin{itemize}
    \item \textbf{CVA (Credit Valuation Adjustment):} The expected loss due to the counterparty defaulting before the trade matures. It is effectively the price of a contingent Credit Default Swap embedded in the derivative. The higher the counterparty's credit risk, the larger the CVA.
    
    \item \textbf{FVA (Funding Valuation Adjustment):} The cost or benefit of funding the collateral (Initial Margin and Variation Margin) required for the trade. If funding is expensive (high rates, poor credit), FVA increases the true cost of the trade.
    
    \item \textbf{KVA (Capital Valuation Adjustment):} The cost of holding regulatory capital (Basel III/IV) over the lifetime of the transaction. Regulators require banks to hold capital as a buffer, and this capital is expensive.
\end{itemize}

Understanding XVA is critical for derivatives pricing in the post-2008 world.

\subsection{Volatility: Market Dynamics and Risk Sensitivity}

\subsubsection{VIX: The Market's Fear Gauge}

The VIX (Volatility Index) measures the market's expectation of 30-day future volatility of the S\&P 500. It is calculated using a range of out-of-the-money (OTM) call and put options, capturing the market's uncertainty about future price moves.

\paragraph{Key Characteristics}
\begin{itemize}
    \item \textbf{Mean-reverting:} The VIX typically hovers around 15-20 in calm markets but spikes to 40-80 during crises, then mean-reverts.
    \item \textbf{Negatively correlated with equity returns:} When stocks fall sharply, the VIX spikes. This ``fear'' relationship is consistent and valuable for portfolio hedging.
    \item \textbf{Leading indicator:} Rising VIX often precedes market downturns, providing an early warning signal.
\end{itemize}

\subsubsection{Greeks: Risk Sensitivities of Options}

Greeks measure the sensitivity of an option's value to changes in underlying factors. The most important Greeks are:

\begin{itemize}
    \item \textbf{Delta ($\Delta$):} Sensitivity to the underlying asset price.
    \item \textbf{Gamma ($\Gamma$):} Sensitivity of delta to changes in the underlying price (second derivative).
    \item \textbf{Theta ($\Theta$):} Time decay—sensitivity to the passage of time.
    \item \textbf{Vega ($\nu$):} Sensitivity to implied volatility.
    \item \textbf{Rho ($\rho$):} Sensitivity to interest rates.
\end{itemize}

\subsubsection{Volga (Vomma): Second-Order Volatility Risk}

Volga measures the sensitivity of Vega to changes in implied volatility—it's the second derivative with respect to volatility:
\begin{equation}
    \text{Volga} = \frac{\partial \text{Vega}}{\partial \sigma} = \frac{\partial^2 V}{\partial \sigma^2}
\end{equation}

\paragraph{Intuition and Practical Importance}

High Volga implies that as volatility increases, your Vega (exposure to volatility) actually grows. This is common in long out-of-the-money (OTM) option positions:

\begin{itemize}
    \item \textbf{At inception:} An OTM option has low Vega because it's unlikely to end up in-the-money.
    \item \textbf{As volatility rises:} The option becomes increasingly valuable, and its Vega increases.
    \item \textbf{Result:} The trader's exposure to further volatility changes becomes amplified.
\end{itemize}

This creates a risk: a trader with high positive Volga is exposed to volatility-of-volatility, which can produce large losses if volatility suddenly drops after rising.